\section{Introduction}
\label{sec:introduction}

Knowledge editing and training data attribution represent two fundamental approaches to understanding factual knowledge in large language models.
Knowledge editing methods such as ROME \citep{meng2022locating} and MEMIT \citep{meng2022mass} perform targeted rank-one updates to MLP weight matrices, treating knowledge as localized key-value mappings in parameter space.
Training data attribution (TDA) methods based on influence functions \citep{koh2017understanding} and gradient-based techniques \citep{park2023trak} estimate which training examples are responsible for a model's factual predictions by computing parameter-space gradients.
Both paradigms operate on the same weight matrices from complementary directions---editing modifies parameters to change knowledge, while attribution measures parameter sensitivity to identify knowledge sources---yet no prior work has examined whether these two notions of ``knowledge direction'' share geometric structure.

This gap has concrete consequences.
\citet{hase2023does} demonstrated a persistent disconnect between knowledge localization (via causal tracing) and knowledge editing success, yet offered no parameter-level explanation.
The editing and attribution communities have developed in isolation, each with independent evaluation paradigms.
A basic geometric question remains unanswered: when ROME identifies a parameter-space direction for editing a fact, and TDA identifies a direction for attributing the same fact to training data, do these directions share any geometric structure?

In this paper, we introduce the TDA-Editing Consistency Score (TECS), a cosine similarity metric between rank-one editing updates and aggregated attribution gradients at the editing layer, and apply it to GPT-2-XL across CounterFact facts with five null baselines and Bonferroni correction.
TECS is indistinguishable from noise (Cohen's $d = 0.050$).
A six-component geometric analysis framework reveals this null result arises from a structured incommensurability under BM25-based attribution: editing directions span a ${\sim}40$-dimensional manifold while attribution gradients collapse to an effectively one-dimensional subspace ($34{:}1$ dimensionality ratio), with asymmetric cross-projection ($17.3\%$ vs.\ $1.0\%$) and principal angles comparable to random baselines.
A theoretical rank-one decomposition predicts that even weak correlations should produce detectable signal under the linear associative memory model, transforming the null result into a quantitative constraint on how severely these assumptions fail.

Our contributions are:
\begin{itemize}
\item We propose TECS, the first metric for comparing knowledge editing directions and attribution gradients in parameter space, with a theoretical rank-one decomposition that yields testable predictions under the linear associative memory model.
\item We characterize structured geometric incommensurability between editing and attribution subspaces---a $34{:}1$ effective dimensionality ratio, large principal angles, and asymmetric cross-projection---providing a parameter-level perspective on the localization-editing disconnect \citep{hase2023does}.
\item We develop a six-component geometric analysis framework and a positive control methodology using a toy linear associative memory to validate metric sensitivity, establishing that the null result in real transformers is informative about knowledge geometry rather than metric failure.
\end{itemize}
